%% mexam-example.tex -- example for the mexam class.
%%
%% Build with latexmk (two runs; see doc/NOTES.md).
%% Class option "tabletitle" swaps the boxed title for the two-by-two cover
%% table built from \exam / \title / \author / \date.
\documentclass[onehalfspacing, 11pt, a4paper]{mexam}

\usepackage{mstuff}
\usepackage{msheet}

\title{Mündliche Prüfung -- Analysis}
\author{Marc Schlienger}
\date{\today}

\exam{Abiturprüfung 2026}      % shown by the "tabletitle" option
\tasknumber{3}                 % shown in the head

\license{\href{https://creativecommons.org/licenses/by/4.0}{CC BY 4.0}}

%%%%%%%%%%%%%%%%%%%%%%%%%%%%%%%%%%%%%%%%%%%%%%%%%%%%%%%%%%%%%%%%%%%%%%%%%%%%%%
\begin{document}

\maketitle

\msection[Ableitung und Monotonie]
Gegeben ist die Funktion $f$ mit $f(x) = \frac{1}{3}x^3 - x$.
\begin{menumerate}
	\item Bestimmen Sie die Nullstellen von $f$. \mscore{3}
	\item Untersuchen Sie $f$ auf Extremstellen. \mscore{4}
	\item Skizzieren Sie den Graphen von $f$ im Intervall $[-3, 3]$. \mscore{3}
\end{menumerate}

\begin{minfo*}
	Erwartungshorizont: Der Prüfling erkennt $f'(x) = x^2 - 1$ und begründet
	das Monotonieverhalten über das Vorzeichen von $f'$.
\end{minfo*}

\msectionr[Impuls für das Prüfungsgespräch]
Wie ändert sich das Schaubild, wenn man $f$ durch $f(x) + c$ ersetzt?
\mscore{2}

\mtotalscore

\end{document}
